% =============================================================================
\bloco{1}{Exam level: the golden rule and the two-period household}{Kurlat \S\S4.3, 6.2; Jones ch.~5 \,\textbullet\, 35 points}

\questao{(8 points) \textit{Source: Kurlat, P. (2020), \emph{A Course in Modern Macroeconomics}, Exercise 5.8 ``National Accounts and
The Golden Rule'', pp.~97--98 --- statement reproduced verbatim, all items (a)--(d).}
\emph{What this trains:} reading $\alpha$, $s$ and $\delta$ off the national accounts, then deciding on which side of the golden rule
the economy sits, by two independent routes (the saving rate against $\alpha$; the MPK against $n+\delta$).}

\noindent In 2018, the following events happened in Aurum, which is a closed economy that uses the Denarius (plural: Denarii) as its currency.
\begin{itemize}[leftmargin=2.2em,topsep=2pt,itemsep=1pt]
\item Grano, Inc.\ hired Cornelia to plant and harvest 500 tons of wheat on a large plot of land using a tractor. Both the land and
the tractor are owned by Grano, Inc.
  \begin{itemize}[label=--,topsep=0pt,itemsep=0pt]
  \item It paid Cornelia 200 Denarii for her work on the harvest.
  \item It sold the wheat to Panem, Inc.\ for 1,000 Denarii.
  \end{itemize}
\item The tractor the Cornelia used:
  \begin{itemize}[label=--,topsep=0pt,itemsep=0pt]
  \item Was new at the beginning of the year. It was worth 2,000 Denarii.
  \item Was no longer new at the end of the year. By then it was worth 1,800 Denarii.
  \end{itemize}
\item Panem, Inc.\ hired Gaius to grind the wheat in its mill and produce a million loaves of bread.
  \begin{itemize}[label=--,topsep=0pt,itemsep=0pt]
  \item It paid Gaius 1,400 Denarii for his work.
  \item It sold the bread for 4,000 Denarii to the hungry citizens, who eat it.
  \end{itemize}
\item Fabrica, Inc.\ hired Livia to build a combine harvester.
  \begin{itemize}[label=--,topsep=0pt,itemsep=0pt]
  \item It paid Livia 400 Denarii for her work.
  \item It sold the combine harvester to Grano, Inc.\ for 1,000 Denarii.
  \end{itemize}
\end{itemize}
\begin{enumerate}[label=(\alph*),leftmargin=2.2em,topsep=2pt,itemsep=1pt]
\item Construct GDP accounts (in Denarii) for Aurum by production, income and expenditure.
\item What were the income shares of labor and capital?
\item What was the savings rate?
\item Assume the following:
  \begin{itemize}[topsep=0pt,itemsep=0pt]
  \item the savings rate is constant,
  \item the production function has the form:
  \[
  Y=AK^\alpha L^{1-\alpha}
  \]
  and there is no technological progress,
  \item the depreciation rate for the tractor is typical for the capital stock as a whole,
  \item the population is constant.
  \end{itemize}
In the long run, would the level of consumption be higher if the economy slightly increased its saving rate?
\end{enumerate}
\resolespaco
\begin{sol}
\begin{notacao}
\item Grano, Panem, Fabrica: the three firms of the exercise. $VA_j$: value added of firm $j$.
\item $C$, $I$, $G$, $NX$: consumption, investment, government purchases and net exports (final uses); here $G=NX=0$.
\item $Y$: GDP. $K$: capital stock. $L$: labour. $A>0$: total factor productivity. $\alpha\in(0,1)$: capital share in
$Y=AK^\alpha L^{1-\alpha}$.
\item $S$: saving. $s\in(0,1)$: saving rate. $\delta\in(0,1)$: depreciation rate. $n\ge0$: population growth rate.
\item $k_{SS}$, $c_{SS}$: steady-state capital and consumption per worker. $s_{GR}$: golden-rule saving rate.
MPK: marginal product of capital, $\partial Y/\partial K$.
\end{notacao}
\passo{(a) production}{The product approach nets out intermediate inputs:}
\begin{equation}
VA_j=\text{sales}_j-\text{intermediate inputs}_j,\qquad \text{GDP}=\textstyle\sum_j VA_j. \label{q1:va}
\end{equation}
The tractor is capital; its wear is depreciation, not an intermediate. The wheat is used up within the year, so it is an
intermediate input of Panem. Using \eqref{q1:va} for each firm:
\begin{align*}
VA_{\text{Grano}} &= 1000-0=1000,\\
VA_{\text{Panem}} &= 4000-1000=3000,\\
VA_{\text{Fabrica}} &= 1000-0=1000.
\end{align*}
Summing with \eqref{q1:va}:
\[
\text{GDP}=1000+3000+1000=5000\ \text{Denarii}.
\]
\passo{(a) income}{GDP is also the sum of factor incomes:}
\begin{equation}
\text{GDP}=\text{wages}+\text{capital income}. \label{q1:inc}
\end{equation}
Wages: $200+1400+400=2000$. Capital income is each firm's value added minus its wages (gross operating surplus, depreciation
included):
\begin{align*}
\text{Grano} &: 1000-200=800,\\
\text{Panem} &: 3000-1400=1600,\\
\text{Fabrica} &: 1000-400=600,\\
\text{total} &: 800+1600+600=3000.
\end{align*}
By \eqref{q1:inc}:
\[
\text{GDP}=2000+3000=5000 .
\]
\passo{(a) expenditure}{GDP is the sum of final uses:}
\begin{equation}
\text{GDP}=C+I+G+NX. \label{q1:exp}
\end{equation}
$C=4000$ (bread eaten by citizens), $I=1000$ (the combine harvester), $G=NX=0$. By \eqref{q1:exp}:
\[
\text{GDP}=4000+1000=5000 .
\]
The three approaches agree.
\passo{(b) shares}{Divide each income total of (a) by GDP:}
\[
\text{labour share}=\frac{2000}{5000}=0.4,\qquad \text{capital share}=\frac{3000}{5000}=0.6 .
\]
With $Y=AK^\alpha L^{1-\alpha}$ and competitive factor markets the capital share is $\alpha$, so $\alpha=0.6$.
\passo{(c) saving rate}{Closed economy, no government: saving is output not consumed, and it equals investment,}
\begin{equation}
S=Y-C=I. \label{q1:S}
\end{equation}
Dividing \eqref{q1:S} by $Y$:
\[
s=\frac{I}{Y}=\frac{1000}{5000}=0.2 .
\]
\passo{(d) route 1, saving rates}{The Solow steady state per worker is}
\begin{align}
k_{SS} &= \left(\frac{sA}{n+\delta}\right)^{\frac1{1-\alpha}}, \label{q1:kss}\\
c_{SS} &= (1-s)Ak_{SS}^\alpha. \label{q1:css}
\end{align}
Substitute \eqref{q1:kss} into \eqref{q1:css}; the exponent becomes $\frac{\alpha}{1-\alpha}$:
\begin{align}
c_{SS} &= (1-s)A\left[\left(\frac{sA}{n+\delta}\right)^{\frac1{1-\alpha}}\right]^{\alpha}\nonumber\\
&= (1-s)A\left(\frac{sA}{n+\delta}\right)^{\frac{\alpha}{1-\alpha}}\nonumber\\
&\propto (1-s)s^{\frac{\alpha}{1-\alpha}}. \label{q1:prop}
\end{align}
With $\alpha=0.6$ from (b), $\frac{\alpha}{1-\alpha}=1.5$, so \eqref{q1:prop} reads
\begin{align}
c_{SS} &\propto (1-s)s^{\frac{(0.6)}{1-(0.6)}}\nonumber\\
&= (1-s)s^{1.5}. \label{q1:c15}
\end{align}
Differentiate \eqref{q1:c15} with the product rule:
\begin{align}
\frac{d}{ds}\big[(1-s)s^{1.5}\big] &= 1.5s^{0.5}(1-s)-s^{1.5}\nonumber\\
&= 1.5s^{0.5}-2.5s^{1.5}\nonumber\\
&= s^{0.5}(1.5-2.5s). \label{q1:deriv}
\end{align}
\eqref{q1:deriv} is zero at $s_{GR}=1.5/2.5=0.6=\alpha$ and positive for $s<0.6$. At $s=0.2$:
\begin{align*}
(0.2)^{0.5}\big(1.5-2.5(0.2)\big) &= 0.2^{0.5}(1.5-0.5)\\
&= 0.447>0 .
\end{align*}
\passo{(d) route 2, MPK}{The tractor, worth 2000 when new, is worth 1800 at year end, and the population is constant:}
\[
\delta=\frac{2000-1800}{2000}=0.1,\qquad n=0 .
\]
In the steady state with $n=0$, saving replaces depreciation:
\begin{equation}
sY=\delta K. \label{q1:ssK}
\end{equation}
Divide \eqref{q1:ssK} by $\delta Y$:
\[
\frac KY=\frac s\delta=\frac{0.2}{0.1}=2 .
\]
Differentiating $Y=AK^\alpha L^{1-\alpha}$ with respect to $K$ gives
\begin{equation}
\text{MPK}=\alpha\frac YK. \label{q1:mpk}
\end{equation}
Substituting $\alpha=0.6$ and $K/Y=2$ into \eqref{q1:mpk}:
\begin{align*}
\text{MPK} &= \frac{(0.6)}{(2)}\\
&= 0.3>n+\delta=0.1 .
\end{align*}
Too little capital: setting \eqref{q1:mpk} equal to $n+\delta=\delta$, the golden rule needs
\begin{align*}
(0.6)\frac YK &= (0.1)\\
\frac KY &= \frac\alpha\delta=6 .
\end{align*}
\passo{answer}{Yes: a slightly higher saving rate raises long-run consumption, because $s=0.2<s_{GR}=\alpha=0.6$ by
\eqref{q1:deriv} (equivalently MPK $=0.3>n+\delta=0.1$ by \eqref{q1:mpk}).}
\intuicao{Below the golden rule an extra unit of capital yields more output (MPK) than it costs to maintain ($n+\delta$), so the
steady state can afford both more capital and more consumption.}
\refbib{Kurlat (2020), \S4.3, eq.\ (4.3.4), p.~61. Matches Problem 5.8 in the project's worked solutions,
\texttt{Resolucao/}\allowbreak\texttt{kurlat\_solutions\_ch05.tex}.}
\rubrica{
\item (a) three approaches, each computed: 3 pts. Wheat treated as final or the harvester netted out: max.\ 1 pt for (a).
\item (b)--(c): 2 pts. (d) both routes: 3 pts; verdict without a computed comparison: max.\ 1 pt for (d).
}
\end{sol}

% -----------------------------------------------------------------------------
\questao{(10 points) \textit{Source: Jones, C.\,I. (2020), \emph{Macroeconomics}, 5th ed., ch.~5, Review Question~4, p.~129 ---
statement reproduced verbatim; item (a) asks it in Jones's model with the exponent $1/3$; items (b)--(d) are added by this list.}
\emph{What this trains:} the steady-state consumption formula with the right exponent, and the exam's exact algebra slip.}

\noindent\textbf{Review Question 4 (Jones).} What is the solution of the Solow model for consumption per person in the steady state,
$c^\ast\equiv C^\ast/L^\ast$? How does each parameter in the solution affect $c^\ast$, and why?

\textbf{a)} (2 points) Answer Review Question 4 in Jones's model ($Y=\bar AK^{1/3}\bar L^{2/3}$, $\Delta K=sY-\bar dK$).

\noindent\textit{Added parts (not in the book):}\par\nopagebreak
\textbf{b)} (2 points) Find the $s$ that maximises $c^\ast$. Check it with the MPK condition.

\textbf{c)} (3 points) In the exam, a student substituted $k^\ast$ instead of $k^{\ast\alpha}$ and maximised $(1-s)\left(\frac{s}{n+\delta}\right)^{\frac1{1-\alpha}}$.
For a general capital share $\alpha$, which saving rate does that give? Compare it with the true golden rule as $\alpha\to0$ and
explain which one passes this knife-edge test, and why.

\textbf{d)} (3 points) Let $Y_t=K_t^\alpha(A_tL_t)^{1-\alpha}$, $A_{t+1}=(1+g)A_t$, $L_{t+1}=(1+n)L_t$, $K_{t+1}=(1-\delta)K_t+sY_t$. With
$\tilde k\equiv K/(AL)$, derive $\tilde k_{t+1}(\tilde k_t)$, the steady state $\tilde k^\ast$, consumption per effective worker $\hat c^\ast(s)$ and
the golden-rule $s$. Show that at any steady state MPK $=\alpha N/s$, with $N\equiv(1+n)(1+g)-1+\delta$. Evaluate $\hat c^\ast$ at
$s\in\{0.2,\ s_{GR},\ 0.6\}$ for $\alpha=1/3$, $n=0.01$, $g=0.02$, $\delta=0.05$.
\resolespaco
\begin{sol}
\begin{notacao}
\item $Y$: output. $K$: capital. $\bar L>0$: labour, constant. $\bar A>0$: TFP, constant. $\Delta K$: change in $K$ over a period.
\item $s\in(0,1)$: saving rate. $\bar d\in(0,1)$: Jones's depreciation rate (Kurlat's $\delta$).
\item $K^\ast$: steady-state capital; $y^\ast\equiv Y^\ast/\bar L$, $c^\ast$: steady-state output and consumption per worker.
\item Items (c)--(d): $\alpha\in(0,1)$ capital share; $n>0$, $g>0$: growth rates of $L$ and $A$; $\delta\in(0,1)$ depreciation.
\item $\tilde k\equiv K/(AL)$, $\tilde y\equiv Y/(AL)$: capital and output per effective worker; $\hat c^\ast$: steady-state consumption
per effective worker. $N\equiv(1+n)(1+g)-1+\delta$: the effective break-even rate. MPK $=\partial Y/\partial K$.
\end{notacao}
\passo{(a) steady state}{Jones's model, with $\bar L$ constant:}
\begin{align}
Y &= \bar AK^{1/3}\bar L^{2/3}, \label{q2:Y}\\
\Delta K &= sY-\bar dK. \label{q2:dK}
\end{align}
Set $\Delta K=0$ in \eqref{q2:dK}:
\begin{equation}
sY^\ast=\bar dK^\ast. \label{q2:ss}
\end{equation}
Substitute \eqref{q2:Y} into \eqref{q2:ss}:
\[
s\big(\bar AK^{\ast1/3}\bar L^{2/3}\big)=\bar dK^\ast .
\]
Divide by $\bar dK^{\ast1/3}$:
\[
K^{\ast2/3}=\frac{s\bar A}{\bar d}\,\bar L^{2/3} .
\]
Raise both sides to the power $3/2$ and divide by $\bar L$:
\begin{equation}
\frac{K^\ast}{\bar L}=\left(\frac{s\bar A}{\bar d}\right)^{3/2}. \label{q2:K}
\end{equation}
Divide \eqref{q2:Y} by $\bar L$ and substitute \eqref{q2:K}:
\begin{align}
y^\ast &= \bar A\left(\frac{K^\ast}{\bar L}\right)^{1/3}\nonumber\\
&= \bar A\left[\left(\frac{s\bar A}{\bar d}\right)^{3/2}\right]^{1/3}\nonumber\\
&= \bar A\left(\frac{s\bar A}{\bar d}\right)^{1/2}\nonumber\\
&= \bar A^{3/2}\left(\frac{s}{\bar d}\right)^{1/2}. \label{q2:ystar}
\end{align}
Consumption is the share $1-s$ of \eqref{q2:ystar}:
\begin{align}
c^\ast &= (1-s)y^\ast\nonumber\\
&= (1-s)\left[\bar A^{3/2}\left(\frac{s}{\bar d}\right)^{1/2}\right]\nonumber\\
&= (1-s)\bar A^{3/2}\left(\frac{s}{\bar d}\right)^{1/2}. \label{q2:cstar}
\end{align}
\passo{(a) parameters}{Read the exponents of \eqref{q2:cstar}. $\partial c^\ast/\partial\bar A>0$ with elasticity $3/2$: TFP raises output
directly and through the capital it induces. $\partial c^\ast/\partial\bar d<0$ (elasticity $-1/2$): more depreciation means less capital.
$s$: ambiguous, a larger capital stock against a smaller consumed share.}
\passo{(b) maximise}{Take the log of \eqref{q2:cstar}:}
\begin{equation}
\ln c^\ast=\text{const}+\ln(1-s)+\tfrac12\ln s. \label{q2:lnc}
\end{equation}
Differentiate \eqref{q2:lnc} with respect to $s$ and set it to zero:
\begin{equation}
-\frac1{1-s}+\frac1{2s}=0. \label{q2:foc}
\end{equation}
Multiply \eqref{q2:foc} by $2s(1-s)>0$:
\begin{align*}
-2s+(1-s) &= 0\\
2s &= 1-s\\
s &= \tfrac13 .
\end{align*}
SOC: differentiating \eqref{q2:foc} once more, $-\frac1{(1-s)^2}-\frac1{2s^2}<0$, a maximum.
Check with the MPK. Differentiating \eqref{q2:Y} and using \eqref{q2:ss}:
\begin{align}
\text{MPK} &= \frac13\,\frac{Y}{K}\nonumber\\
&= \frac13\left(\frac{\bar d}{s}\right). \label{q2:mpk}
\end{align}
Substituting $s=1/3$ into \eqref{q2:mpk}:
\begin{align*}
\text{MPK} &= \frac13\left(\frac{\bar d}{(1/3)}\right)\\
&= \bar d,
\end{align*}
so MPK equals $\bar d$ (here $n=0$) exactly at $s=1/3$.
\passo{(c) the slip}{The student's objective, in logs:}
\begin{equation}
\ln\big[(1-s)s^{\frac1{1-\alpha}}\big]=\ln(1-s)+\frac{1}{1-\alpha}\ln s. \label{q2:slip}
\end{equation}
Differentiate \eqref{q2:slip} and set it to zero:
\[
-\frac1{1-s}+\frac{1}{(1-\alpha)s}=0 .
\]
Multiply by $(1-\alpha)s(1-s)>0$:
\begin{align}
(1-\alpha)s &= 1-s\nonumber\\
s &= \frac1{2-\alpha}. \label{q2:sslip}
\end{align}
The right objective carries the exponent $\frac{\alpha}{1-\alpha}$:
\begin{equation}
\ln\big[(1-s)s^{\frac{\alpha}{1-\alpha}}\big]=\ln(1-s)+\frac{\alpha}{1-\alpha}\ln s. \label{q2:right}
\end{equation}
Differentiate \eqref{q2:right}, set it to zero and multiply by $(1-\alpha)s(1-s)$:
\begin{align}
\alpha(1-s) &= (1-\alpha)s\nonumber\\
s_{GR} &= \alpha. \label{q2:sgr}
\end{align}
At $\alpha=1/3$, \eqref{q2:sslip} gives
\begin{align*}
s &= \frac1{2-(1/3)}\\
&= 0.6,
\end{align*}
against $s_{GR}=1/3$ from \eqref{q2:sgr}.
\passo{(c) knife edge}{As $\alpha\to0$ capital produces nothing, so every unit saved is wasted and the best $s$ must go to $0$:
\eqref{q2:sgr} gives $s_{GR}=\alpha\to0$ and passes; \eqref{q2:sslip} gives $\frac1{2-\alpha}\to\frac12$ and fails (it would save half of
output to maintain useless machines).}
\passo{(d) law of motion}{The statement's model:}
\begin{align}
Y_t &= K_t^\alpha(A_tL_t)^{1-\alpha}, \label{q2:Yd}\\
K_{t+1} &= (1-\delta)K_t+sY_t, \label{q2:Kd}\\
A_{t+1}L_{t+1} &= (1+g)(1+n)A_tL_t. \label{q2:AL}
\end{align}
Divide \eqref{q2:Yd} by $A_tL_t$:
\begin{equation}
\tilde y_t=\tilde k_t^\alpha. \label{q2:yt}
\end{equation}
Divide \eqref{q2:Kd} by $A_{t+1}L_{t+1}$, written as in \eqref{q2:AL}, and use \eqref{q2:yt}:
\begin{align}
\tilde k_{t+1} &= \frac{(1-\delta)K_t+sY_t}{\big((1+g)(1+n)A_tL_t\big)}\nonumber\\
&= \frac{(1-\delta)\tilde k_t+s\tilde y_t}{(1+n)(1+g)}\nonumber\\
&= \frac{(1-\delta)\tilde k_t+s\big(\tilde k_t^\alpha\big)}{(1+n)(1+g)}\nonumber\\
&= \frac{(1-\delta)\tilde k_t+s\tilde k_t^\alpha}{(1+n)(1+g)}. \label{q2:lom}
\end{align}
Steady state: set $\tilde k_{t+1}=\tilde k_t=\tilde k$ in \eqref{q2:lom} and multiply by $(1+n)(1+g)$:
\begin{align}
(\tilde k) &= \frac{(1-\delta)(\tilde k)+s(\tilde k)^\alpha}{(1+n)(1+g)}\nonumber\\
(1+n)(1+g)\tilde k &= (1-\delta)\tilde k+s\tilde k^\alpha\nonumber\\
\big[(1+n)(1+g)-1+\delta\big]\tilde k &= s\tilde k^\alpha\nonumber\\
N\tilde k &= s\tilde k^\alpha. \label{q2:ssd}
\end{align}
Solve \eqref{q2:ssd} for $\tilde k$:
\begin{equation}
\tilde k^\ast=\left(\frac sN\right)^{\frac1{1-\alpha}}. \label{q2:kt}
\end{equation}
Consumption per effective worker is $(1-s)\tilde y$; use \eqref{q2:yt} and \eqref{q2:kt}:
\begin{align}
\hat c^\ast &= (1-s)\big(\tilde k^{\ast\alpha}\big)\nonumber\\
&= (1-s)\left[\left(\frac sN\right)^{\frac1{1-\alpha}}\right]^{\alpha}\nonumber\\
&= (1-s)\left(\frac sN\right)^{\frac{\alpha}{1-\alpha}}. \label{q2:chat}
\end{align}
In $s$, \eqref{q2:chat} is the objective \eqref{q2:right} up to a constant, so it is maximised at $s_{GR}=\alpha$, \eqref{q2:sgr}.
Expanding, $N=n+g+ng+\delta$: the continuous-time $n+g+\delta$ plus the cross term $ng$.
\passo{(d) MPK cross-check}{Differentiate \eqref{q2:Yd} with respect to $K_t$:}
\begin{equation}
\text{MPK}=\alpha K^{\alpha-1}(AL)^{1-\alpha}=\alpha\tilde k^{\alpha-1}. \label{q2:mpkd}
\end{equation}
By \eqref{q2:kt}, $\tilde k^{\ast1-\alpha}=s/N$. Substitute into \eqref{q2:mpkd}:
\begin{align}
\text{MPK} &= \frac{\alpha}{\tilde k^{\ast1-\alpha}}\nonumber\\
&= \frac{\alpha}{(s/N)}\nonumber\\
&= \frac{\alpha N}{s}. \label{q2:mpkN}
\end{align}
So MPK $=N\iff s=\alpha$; $s<\alpha\iff$ MPK $>N$ (too little capital).
\passo{(d) numbers}{The break-even rate:}
\[
N=1.01\times1.02-1+0.05=0.0802 .
\]
With $\alpha=1/3$, $\frac1{1-\alpha}=1.5$. Using \eqref{q2:kt}:
\begin{align*}
s=0.2 &: \ \tilde k^\ast=(0.2/0.0802)^{1.5}=3.9381,\\
s=1/3 &: \ \tilde k^\ast=(0.3333/0.0802)^{1.5}=8.4734,\\
s=0.6 &: \ \tilde k^\ast=(0.6/0.0802)^{1.5}=20.4628.
\end{align*}
Using \eqref{q2:chat} as $(1-s)\tilde k^{\ast1/3}$:
\begin{align*}
s=0.2 &: \ \hat c^\ast=0.8\times3.9381^{1/3}=1.2633,\\
s=1/3 &: \ \hat c^\ast=\tfrac23\times8.4734^{1/3}=1.3591,\\
s=0.6 &: \ \hat c^\ast=0.4\times20.4628^{1/3}=1.0941.
\end{align*}
Using \eqref{q2:mpkN}:
\begin{align*}
s=0.2 &: \ \text{MPK}=\frac{(1/3)(0.0802)}{(0.2)}=0.1337,\\
s=1/3 &: \ \text{MPK}=\frac{(1/3)(0.0802)}{(1/3)}=0.0802=N,\\
s=0.6 &: \ \text{MPK}=\frac{(1/3)(0.0802)}{(0.6)}=0.0446.
\end{align*}
The slip's answer
($s=0.6$, from \eqref{q2:sslip}) loses
\[
1-\frac{1.0941}{1.3591}=19.5\%
\]
of the attainable consumption.
\intuicao{The golden rule reinvests exactly capital's income: saving $\alpha$ of output pays the maintenance $N\tilde k$ of the stock whose
marginal product equals that maintenance cost.}
\refbib{Jones (2020), ch.~5, eqs.\ (5.7)--(5.9), pp.~112--113; Jones, \emph{Instructor's Manual} (2016), ch.~5, answer to Review
Question 4, p.~37 (gives $c^\ast=(1-s)\bar A^{3/2}(s/\bar d)^{1/2}$ and ``optimal $s$ = the exponent on capital, 1/3''; its line that a higher
depreciation rate \emph{raises} $c^\ast$ is a slip: $\bar d$ is in the denominator). Kurlat (2020), \S4.3; Romer (2012), Problem 2.15(c), p.~97.}
\rubrica{
\item (a) $c^\ast$ with exponents $3/2$, $1/2$: 2 pts. (b): 2 pts; no MPK check: $-0.5$.
\item (c) $1/(2-\alpha)$ and the knife-edge argument: 3 pts. (d): 3 pts; arithmetic slip with correct method: $-0.5$ (does not add to a ceiling).
\item Ceiling: maximising $(1-s)k^\ast$ instead of $(1-s)k^{\ast\alpha}$ anywhere: max.\ 1 pt on that sub-item.
}
\end{sol}

% -----------------------------------------------------------------------------
\questao{(10 points) \textit{Source: Kurlat, P. (2020), \emph{A Course in Modern Macroeconomics}, Exercise 6.1 ``Two-Period Problem
with Taxes and Initial Wealth'', pp.~121--122 --- statement reproduced verbatim, all items (a)--(e).}
\emph{What this trains:} the full household problem with every income source, closed forms, and how $\sigma$ decides
whether $c_1$ rises or falls with $r$ (the log case $\sigma=1$ is the exam's).}

\noindent Suppose a household solves the following two-period consumption-savings problem with taxes:
\begin{gather*}
\max_{c_1,a,c_2}\ u(c_1)+\beta u(c_2)\\
\text{s.t.}\\
a=a_0+y_1-\tau_1-c_1\\
c_2=y_2-\tau_2+(1+r)\,a
\end{gather*}
with $u(c)=\frac{c^{1-\sigma}}{1-\sigma}$, where: $c_1$ is consumption at time 1, $c_2$ is consumption at time 2, $y_1$ is household income
at time 1, $y_2$ is household income at time 2, $\tau_1$ are taxes at time 1, $\tau_2$ are taxes at time 2, and $a_0$ is initial wealth.
\begin{enumerate}[label=(\alph*),leftmargin=2.2em,topsep=2pt,itemsep=1pt]
\item Solve for the household's choice of $c_1$, $c_2$ and $a$ in closed form.\footnote{Kurlat's footnote: ``In closed form'' means
that you have an explicit expression for something. Suppose I ask you to determine some variable $x$. Then something like
$x=y^2-z$ is a solution ``in closed form''; something like $e^x+x-b=0$ also implicitly tells what $x$ should be but is not in closed form.}
\item How does $\frac{c_1}{y_1}$ depend on $y_2$? What would happen if households suddenly became optimistic about the future?
\item How does $\frac{c_1}{y_1}$ depend on $a_0$? Interpret.
\item How does $\frac{c_1}{y_1}$ depend on $\beta$? Interpret.
\item Suppose $y_2=\tau_1=\tau_2=0$ and compute $\frac{\partial c_1}{\partial r}$. How does the answer depend on $\sigma$? Interpret the
answer. [Hint: this is a hard question, not the maths but the interpretation. Think about what $\sigma$ means for the relative
importance of income and substitution effects]
\end{enumerate}
\resolespaco
\begin{sol}
\begin{notacao}
\item $c_1,c_2>0$: consumption in periods 1 and 2. $a$: assets carried into period 2 ($a<0$ is borrowing).
\item $a_0\ge0$: initial wealth. $y_1,y_2\ge0$: income. $\tau_1,\tau_2$: lump-sum taxes. $r>-1$: interest rate.
\item $\beta\in(0,1)$: discount factor. $u(c)=\frac{c^{1-\sigma}}{1-\sigma}$, $\sigma>0$ (curvature; $1/\sigma$ is the intertemporal
elasticity of substitution; $\sigma=1$ is $\ln c$), so $u'(c)=c^{-\sigma}$.
\item $W$: present value of resources; $D$: the consumption divisor (both defined below, where first used).
\end{notacao}
\passo{(a) set-up}{Choice variables $c_1,c_2,a$. The two period budgets:}
\begin{align}
c_1+a &= a_0+y_1-\tau_1, \label{q3:b1}\\
c_2 &= y_2-\tau_2+(1+r)a. \label{q3:b2}
\end{align}
Solve \eqref{q3:b1} for $a$ and substitute into \eqref{q3:b2}:
\begin{equation}
c_2=y_2-\tau_2+(1+r)(a_0+y_1-\tau_1-c_1). \label{q3:c2}
\end{equation}
With $c_2$ given by \eqref{q3:c2}, the problem has one choice:
\begin{equation}
\max_{c_1}\ u(c_1)+\beta u(c_2),\qquad c_2\ \text{from \eqref{q3:c2}}. \label{q3:max}
\end{equation}
\passo{(a) Euler}{Differentiate \eqref{q3:max} with respect to $c_1$, using $\partial c_2/\partial c_1=-(1+r)$ from \eqref{q3:c2}:}
\begin{equation}
u'(c_1)-\beta(1+r)u'(c_2)=0. \label{q3:foc}
\end{equation}
Substitute $u'(c)=c^{-\sigma}$ into \eqref{q3:foc}:
\begin{align*}
\big(c_1^{-\sigma}\big)-\beta(1+r)\big(c_2^{-\sigma}\big) &= 0\\
c_1^{-\sigma} &= \beta(1+r)c_2^{-\sigma}.
\end{align*}
Divide by $c_1^{-\sigma}\beta(1+r)$ and raise to the power $1/\sigma$:
\begin{equation}
c_2=[\beta(1+r)]^{1/\sigma}c_1. \label{q3:euler}
\end{equation}
\passo{(a) closed forms}{Divide \eqref{q3:c2} by $1+r$ and move $c_1$ to the left: the present-value budget is}
\begin{equation}
c_1+\frac{c_2}{1+r}=W,\qquad W\equiv a_0+y_1-\tau_1+\frac{y_2-\tau_2}{1+r}. \label{q3:pv}
\end{equation}
Substitute \eqref{q3:euler} into \eqref{q3:pv}:
\begin{align*}
c_1+\frac{\big([\beta(1+r)]^{1/\sigma}c_1\big)}{1+r} &= W\\
c_1+\frac{\beta^{1/\sigma}(1+r)^{1/\sigma}}{1+r}\,c_1 &= W\\
c_1\big[1+\beta^{1/\sigma}(1+r)^{\frac1\sigma-1}\big] &= W.
\end{align*}
Name the bracket
\begin{equation}
D\equiv1+\beta^{1/\sigma}(1+r)^{\frac1\sigma-1}>1. \label{q3:D}
\end{equation}
Then
\begin{equation}
c_1=\frac WD. \label{q3:c1}
\end{equation}
Substitute \eqref{q3:c1} into \eqref{q3:euler}, and into \eqref{q3:b1} solved for $a$:
\begin{align*}
c_2 &= [\beta(1+r)]^{1/\sigma}\left(\frac WD\right)\\
&= \frac{[\beta(1+r)]^{1/\sigma}W}{D},\\
a &= a_0+y_1-\tau_1-\left(\frac WD\right)\\
&= a_0+y_1-\tau_1-\frac WD .
\end{align*}
\passo{(b) future income}{By \eqref{q3:pv}, $\partial W/\partial y_2=\frac1{1+r}$. Differentiate \eqref{q3:c1}:}
\begin{align*}
\frac{\partial c_1}{\partial y_2} &= \frac1D\left(\frac1{1+r}\right)\\
&= \frac1{(1+r)D}>0,\\
\frac{\partial(c_1/y_1)}{\partial y_2} &= \frac1{y_1(1+r)D}.
\end{align*}
Optimism about $y_2$ raises consumption and cuts saving today.
\passo{(c) wealth}{By \eqref{q3:pv}, $\partial W/\partial a_0=\partial W/\partial y_1=1$. Differentiate \eqref{q3:c1}:}
\begin{align*}
\frac{\partial c_1}{\partial a_0}=\frac{\partial c_1}{\partial y_1} &= \frac1D\,(1)\\
&= \frac1D\in(0,1),
\end{align*}
using $D>1$ from \eqref{q3:D}. Only the present value of resources matters, and the household spreads any windfall over both periods.
\passo{(d) patience}{Differentiate \eqref{q3:D} with respect to $\beta$:}
\[
\frac{\partial D}{\partial\beta}=\frac1\sigma\beta^{\frac1\sigma-1}(1+r)^{\frac1\sigma-1}>0 .
\]
Differentiate \eqref{q3:c1} ($W$ does not contain $\beta$) and substitute:
\[
\frac{\partial c_1}{\partial\beta}=-\frac{W}{D^2}\cdot\frac{\partial D}{\partial\beta}
=-\frac{W}{D^2}\cdot\left(\frac1\sigma\beta^{\frac1\sigma-1}(1+r)^{\frac1\sigma-1}\right)<0 .
\]
A more patient household consumes less now and saves more.
\passo{(e) interest rate}{With $y_2=\tau_1=\tau_2=0$, \eqref{q3:pv} gives}
\begin{align*}
W &= a_0+y_1-(0)+\frac{(0)-(0)}{1+r}\\
&= a_0+y_1,
\end{align*}
so \eqref{q3:c1} becomes
\begin{equation}
c_1=\frac{(a_0+y_1)}D. \label{q3:c1e}
\end{equation}
Differentiate \eqref{q3:D} with respect to $r$:
\begin{equation}
D'(r)=\beta^{1/\sigma}\left(\frac1\sigma-1\right)(1+r)^{\frac1\sigma-2}. \label{q3:Dr}
\end{equation}
Differentiate \eqref{q3:c1e} and substitute \eqref{q3:Dr}, using $-\left(\frac1\sigma-1\right)=\frac{\sigma-1}{\sigma}$:
\begin{align}
\frac{\partial c_1}{\partial r} &= -\frac{(a_0+y_1)D'(r)}{D^2}\nonumber\\
&= -\frac{(a_0+y_1)\left(\beta^{1/\sigma}\left(\frac1\sigma-1\right)(1+r)^{\frac1\sigma-2}\right)}{D^2}\nonumber\\
&= \frac{(a_0+y_1)\beta^{1/\sigma}(1+r)^{\frac1\sigma-2}}{D^2}\cdot\frac{\sigma-1}{\sigma}. \label{q3:dc1dr}
\end{align}
Sign of \eqref{q3:dc1dr} $=$ sign of $\sigma-1$. A saver is richer when $r$ rises (income effect: $c_1\uparrow$), and waiting is cheaper
(substitution: $c_1\downarrow$). $\sigma>1$ (low elasticity $1/\sigma$): income wins. $\sigma<1$: substitution wins. $\sigma=1$: they
cancel; \eqref{q3:D} gives
\begin{align*}
D &= 1+\beta^{1/(1)}(1+r)^{\frac1{(1)}-1}\\
&= 1+\beta,
\end{align*}
and $c_1=\frac{a_0+y_1}{1+\beta}$ does not depend on $r$.
\intuicao{Numerator $W$ = how much the household has; denominator $D$ = how it splits it across dates. Every comparative static moves one
of the two.}
\refbib{Kurlat (2020), \S6.2, pp.~105--117. Instructor key PSET~3 (\texttt{Listas/soluções do instrutor/PSET 3 (1).pdf}), Q1(a)--(c):
same method, same $c_1$ and $a$. Two slips in that key: its boxed $c_2$ has denominator $1+\beta^{1/\sigma}(1+r)^{\frac1\sigma-1}$ where
$1+\beta^{-1/\sigma}(1+r)^{1-\frac1\sigma}$ is needed, and its $\partial c_1/\partial r$ drops the factor $\beta^{1/\sigma}$; the sign, set by $\frac{\sigma-1}{\sigma}$, is right.}
\rubrica{
\item (a) Euler and closed forms: 4 pts. Euler only, no closed forms: max.\ 2 pts.
\item (b)--(d): 3 pts. (e) derivative and its sign: 2 pts; income and substitution reading: 1 pt.
\item Ceiling: $1/\sigma$ and $\sigma$ swapped in the reading of (e): max.\ 1 pt for (e).
}
\end{sol}

% -----------------------------------------------------------------------------
\questao{(7 points) \textit{Source: Kurlat, P. (2020), \emph{A Course in Modern Macroeconomics}, Exercise 6.2 ``Savings Rates for
Different People'', p.~122 --- statement reproduced verbatim, all items (a)--(b). Its Exercise~6.1 is Question~3 of this list.}
\emph{What this trains:} a saving rate out of \emph{current} income depends on the income profile and, through human wealth, on $r$,
even when the saving rate out of \emph{lifetime} wealth does not.}

\noindent An economist has data on the occupation, this year's income (denoted $y$) and this year's consumption (denoted $c$) of a sample
of 26-year-olds. Within this sample, some are top professional athletes and others are medical doctors in their first year of
residence. [Make whatever assumptions you think are reasonable and, if you wish, refer back to Exercise 6.1]
\begin{enumerate}[label=(\alph*),leftmargin=2.2em,topsep=2pt,itemsep=1pt]
\item Suppose one computed $s=\frac{y-c}{y}$ for each individual in the sample. Should we expect $s$ to be higher for athletes or doctors?
\item Suppose interest rates go down. How should we expect the response of $s$ to differ between the two groups?
\end{enumerate}
\resolespaco
\begin{sol}
\begin{notacao}
\item $y_1,y_2>0$: income in periods 1 and 2; $m\equiv y_2/y_1$: the income profile. $c_1$: consumption in period 1.
\item $s\equiv1-c_1/y_1$: saving rate out of \emph{current} income ($s<0$ is borrowing).
\item $r>-1$, $\beta\in(0,1)$, $\sigma>0$, $W$ and $D$: as in Question~3, \eqref{q3:pv} and \eqref{q3:D}.
\end{notacao}
\passo{(a) assumption}{A 26-year-old athlete earns most of his lifetime income now ($m$ small); a first-year resident earns
little now and much later ($m$ large). Use Question~3 with $a_0=\tau_1=\tau_2=0$; then \eqref{q3:pv} gives}
\begin{equation}
W=y_1+\frac{y_2}{1+r}=y_1\Big(1+\frac{m}{1+r}\Big). \label{q4:W}
\end{equation}
The saving rate out of current income is
\begin{equation}
s=1-\frac{c_1}{y_1}. \label{q4:s}
\end{equation}
Substitute \eqref{q3:c1} and then \eqref{q4:W} into \eqref{q4:s}:
\begin{align}
s &= 1-\frac1{y_1}\left(\frac WD\right)\nonumber\\
&= 1-\frac1{y_1D}\left(y_1\Big(1+\frac{m}{1+r}\Big)\right)\nonumber\\
&= 1-\frac1D\Big(1+\frac{m}{1+r}\Big), \label{q4:sm}
\end{align}
which is decreasing in $m$.
\passo{(a) answer}{$s$ is higher for athletes; residents' $s$ can be negative (they borrow). With $\beta=0.97$, $\sigma=2$, $r=0.06$,
\eqref{q3:D} gives $D=1+0.97^{1/2}\,1.06^{-1/2}=1.9566$. Using \eqref{q4:sm}:}
\begin{align*}
\text{athlete } (y_1,y_2)=(10,2),\ m=0.2 &: \ s=1-\frac{1+0.2/1.06}{1.9566}=0.392,\\
\text{resident } (3,9),\ m=3 &: \ s=1-\frac{1+3/1.06}{1.9566}=-0.958.
\end{align*}
\passo{(b) channels}{Differentiate \eqref{q4:W} with respect to $r$:}
\[
\frac{\partial W}{\partial r}=-\frac{y_2}{(1+r)^2}<0 .
\]
So a fall in $r$ (i) raises human wealth $\frac{y_2}{1+r}$, by an amount proportional to $y_2$, raising $c_1$; (ii) changes $D$, whose
sign depends on $\sigma$ (Question~3(e), \eqref{q3:Dr}). For borrowers (residents) income and substitution effects both raise $c_1$;
for lenders (athletes) they conflict.
\passo{(b) answer}{Residents' $s$ falls clearly; athletes' $s$ barely moves (and, with $\sigma>1$, rises). Cut $r$ from $6\%$ to $2\%$;
\eqref{q3:D} and then \eqref{q4:sm} give:}
\begin{align*}
D &= 1+0.97^{1/2}\,(1.02)^{-1/2}=1.9752,\\
\text{athlete} &: \ s=1-\frac{1+0.2/(1.02)}{(1.9752)}=0.3944,\\
\text{resident} &: \ s=1-\frac{1+3/(1.02)}{(1.9752)}=-0.9953.
\end{align*}
Subtracting the values of (a) to four decimals ($0.3925$ and $-0.9576$):
\[
\text{athlete } \Delta s=+0.002,\qquad \text{resident } \Delta s=-0.038 .
\]
\intuicao{The same optimising rule produces opposite saving rates when income arrives at different dates; a rate cut is a transfer from
lenders to borrowers.}
\refbib{Kurlat (2020), \S6.2; numbers as in the project's worked solution \texttt{Resolucao/}\allowbreak\texttt{kurlat\_solutions\_ch06.tex}, Problem 6.2.}
\rubrica{
\item (a) the income-profile argument with the formula for $s$: 3 pts; ranking without the profile argument: max.\ 1 pt.
\item (b) human-wealth channel and borrower/lender sign: 4 pts; ``lower $r$ always lowers saving'' for both groups: max.\ 1 pt.
}
\end{sol}
